\section{Test Approach and Special Considerations}\label{sec:approach}

\subsection{Approach}
The test cases are black-box tests derived from the requirements. For each requirement there is a case for the behavior it specifies and a case for each refusal, limit, or boundary it states; the cases for secondary rules and refusals have Medium priority, and the short edition of this document leaves them out. A case is run by hand through the web client, as the user named in its Actor field, so it shows what that user sees. Where the client has no screen for a check, such as calling an endpoint without the needed role, the step calls the API and says so. Where an automated test runs the same case, the Automated by field names it: those tests run on PostgreSQL~16 in Testcontainers with the production migrations through \code{mvn verify}, and each is tagged with the keys of the cases it runs (\code{@Verifies}), so that \code{mvn verify} also writes a run report, \code{target/test-cases/report.md}, with every case's outcome in that build: Passed, Passed in part, Failed, or Not automated. The cases are run in the order of Section~\ref{sec:cases}, but each case's preconditions are complete, so a single case can be run alone after a reset.

\subsection{Pass and Fail Criteria}
A test case \textbf{passes} when the expected result of every step is observed and its pass criteria hold. ``Observed'' means seen on the screen the step names, or, for a step that calls the API, in the response's status and body. The Request detail's history, the Request list, the Inventory, and the Team's Statement are the places where state, Holds, and Transactions are checked. A test case \textbf{fails} when any expected result is not observed, or when the system changes anything the case says must not change, such as a Statement that gains a Transaction after a refused submission. A case that cannot be run because an earlier case or a precondition failed is \textbf{blocked}, not failed. Each run records the date, the build, the tester, the outcome of each case, and for a failure the step, what was observed, and the defect raised in the team's issue tracker.

\subsection{Entry and Exit Criteria}
\textbf{Entry.} The build under test compiles, \code{mvn verify} passes, and the application starts with the demo profile and shows the test world of Section~\ref{sec:world}. \textbf{Exit.} Every test case of a built feature has been run on the build under test and its result recorded; every High priority case passes; every failing case has a defect raised and assigned.

\subsection{Special Considerations}\label{sec:special}
\begin{itemize}
  \item \textbf{Features not yet built.} The cases for features that the build under test does not implement yet (Section~\ref{sec:product}) are written against the SRS and the screens of the SDD, and their Automated by field reads ``None yet''. Until the feature is built they are recorded as \emph{not run}, never as failed, and they are run when the feature's increment is delivered. A case's steps do not change when the feature is built, unless a screen named in the case turns out differently, in which case the step is reworded and the change entered in the Record of Changes. A case that is built in part, such as one whose last step checks a notification or that adds a Person line beside Exclusive and Metered ones, is run up to the first step that needs an unbuilt feature; the steps from there are recorded as not run, and the case counts as passed only once every step has run and passed.
  \item \textbf{Stretch requirements.} The six Stretch requirements have Low priority cases. They are run only if the feature is built for Release~1; otherwise they are recorded as not in the release.
  \item \textbf{Time.} Cases whose steps are marked [clock] depend on time passing, as described in Section~\ref{sec:dates}. Windows must start in the future, so every case uses the relative days of that section.
  \item \textbf{Concurrency.} Two Managers racing for the last room or the last GPU units cannot be timed reliably by hand. Those cases are run by their automated tests, which start both submissions at once; a manual run checks only that the outcome is one allocation and one waiting line.
  \item \textbf{Roles and scoping.} The server enforces roles and Team scope, not the client. A Request, Account, or Statement of another Team answers ``not found'' (\code{NOT\_FOUND}), and an action without its role answers \code{FORBIDDEN}. Cases that check this call the API directly, because the client never offers the forbidden action.
  \item \textbf{Notifications.} Notifications are in-app only. A step that expects a notification checks the recipient's notification list after signing in as that Member.
  \item \textbf{Data.} The test world uses only invented names and the \code{openbench.example} mail domain; no real personal data is used.
\end{itemize}
